\slcol{ಶ್ರೀ ಭಗವಾನುವಾಚ ।\\
\Index{ಮಯ್ಯಾಸಕ್ತಮನಾಃ ಪಾರ್ಥ} ಯೋಗಂ ಯುಂಜನ್ಮದಾಶ್ರಯಃ ।\\
ಅಸಂಶಯಂ ಸಮಗ್ರಂ ಮಾಂ ಯಥಾ ಜ್ಞಾಸ್ಯಸಿ ತಚ್ಛೃಣು ॥ ೧ ॥}
\cquote{ಶ್ರೀ ಭಗವಂತನು  ಹೇಳಿದನು,
ಪಾರ್ಥ, ನನ್ನಲ್ಲಿಯೇ ಮನವನ್ನು ನೆಲೆಸಿ ನನ್ನನ್ನೇ ಶರಣೆಂದು  ನಂಬಿ ಯೋಗ ಸಾಧನವನ್ನು ನಡೆಸುತ್ತಾ ನಿಸ್ಸಂಶಯವಾಗಿ, ಪೂರ್ಣವಾಗಿ ನನ್ನನ್ನು ಹೇಗೆ ತಿಳಿಯಲಾಗುವುದೋ ಅದನ್ನು ಕೇಳು.}
\slcol{\Index{ಜ್ಞಾನಂ ತೇಽಹಂ} ಸವಿಜ್ಞಾನಮಿದಂ ವಕ್ಷ್ಯಾಮ್ಯಶೇಷತಃ ।\\
ಯಜ್ಜ್ಞಾತ್ವಾ ನೇಹ ಭೂಯೋಽನ್ಯಜ್ಜ್ಞಾತವ್ಯಮವಶಿಷ್ಯತೇ ॥ ೨ ॥}
\cquote{ಯಾವುದನ್ನು ತಿಳಿದಮೇಲೆ ಈ ಜಗತ್ತಿನಲ್ಲಿ ತಿಳಿಯತಕ್ಕದ್ದು ಎನ್ನುವುದು ಇನ್ನೊಂದು ಉಳಿಯುವುದಿಲ್ಲವೋ, ಅಂತಹ ಸೂಕ್ಷ್ಮ ವಿವರದಿಂದ ಕೂಡಿದ ಈ ಜ್ಞಾನಮಾರ್ಗವನ್ನು ಪೂರ್ಣವಾಗಿ ನಿನಗೆ ನಾನು ಹೇಳುತ್ತೇನೆ.}

\newpage
\begin{mananam}{\mananamTitle{ಮನನ ಶ್ಲೋಕ - ೨}}
\nfntmtxt ನಾನು ಆಚರಣೆಯಲ್ಲಿ ಆಧ್ಯಾತ್ಮಿಕತೆಯನ್ನು ತೊಡಗಿಸಿಕೊಳ್ಳದೆ, ಆಧ್ಯಾತ್ಮಿಕ ಜ್ಞಾನವನ್ನು ಪಡೆಯುವುದರ ಕಡೆಗೆ ಬಹಳವಾಗಿ ಒಲವು ತೋರಿಸುತ್ತಿದ್ದೇನೆಯೆ? ಜೀವನದ ಇತರ ಕ್ಷೇತ್ರಗಳಲ್ಲಿರುವಂತೆ, ಆಧ್ಯಾತ್ಮಿಕ ಜೀವನದಲ್ಲಿ ಸೈದ್ಧಾಂತಿಕ ಜ್ಞಾನಕ್ಕೂ, ಪ್ರಾಯೋಗಿಕ ತಿಳುವಳಿಕೆಗೂ ಇರುವ ವ್ಯತ್ಯಾಸವೇನು? ನನ್ನ ಆಧ್ಯಾತ್ಮಿಕ ಒಳನೋಟಗಳನ್ನು ನನ್ನ ದೈನಂದಿನ ಜೀವನದಲ್ಲಿ ಉಪಯುಕ್ತವಾಗುವಂತೆ, ವಿವೇಕ ಹಾಗೂ ವೈಯಕ್ತಿಕ ಸತ್ಯ ಸ್ಥಿತಿಯನ್ನು ತಿಳಿಯುವುದಕ್ಕೆ  ಹೇಗೆ ರೂಪಾಂತರಿಸಬಹುದು? 
\end{mananam}
\sReflect
\begin{inspiration}{\mananamTitle{ಸ್ಪೂರ್ತಿ}}
\nfntmtxt ಜೀವನದ ಪ್ರತಿಯೊಂದು ಅಂಶಗಳಲ್ಲೂ, ಸೈದ್ಧಾಂತಿಕ ಮತ್ತು ಆಚರಣೆಯ ಜ್ಞಾನದ ನಡುವೆ ವ್ಯತ್ಯಾಸಗಳಿರುತ್ತವೆ. ಆಧ್ಯಾತ್ಮಿಕ ಜ್ಞಾನವು ಎಲ್ಲರಿಗೂ ಸಾಮಾನ್ಯವಾಗಿ ಅನ್ವಯಿಸುವುದಾದರೂ, ಪ್ರತಿಯೊಬ್ಬ ವ್ಯಕ್ತಿಯೂ ಅದನ್ನು ತನ್ನ ಜೀವನದಲ್ಲಿ ಅಳವಡಿಸಿಕೊಳ್ಳುವ ರೀತಿಯಲ್ಲಿ ವ್ಯತ್ಯಾಸವಿರುತ್ತದೆ. ಆದಾಗ್ಯೂ, ಆಧ್ಯಾತ್ಮಿಕ ಜ್ಞಾನದಲ್ಲಿ ಒಂದು ಪ್ರಮುಖ ವ್ಯತ್ಯಾಸವಿದೆ; ಇತರ ಕ್ಷೇತ್ರಗಳಲ್ಲಿ ಕೊನೆಯಿಲ್ಲದ ಅಪಾರವಾದ ಜ್ಞಾನದ ಶೇಖರಣೆಗೆ ಎಡೆ ಇರುತ್ತದೆ ಆದರೆ, ಇದಕ್ಕೆ ವ್ಯತಿರಿಕ್ತವಾಗಿ ಆಧ್ಯಾತ್ಮಿಕ ಜ್ಞಾನಕ್ಕೆ ಒಂದು ಕೊನೆ ಎಂಬುದಿರುತ್ತದೆ.  ಯಾವ ಒಂದೇ ಒಂದು ಜ್ಞಾನ ತಿಳಿಯುವುದರಿಂದ ಬೇರೆಲ್ಲಾ ಜ್ಞಾನಗಳೂ ತಮ್ಮಷ್ಟಕ್ಕೆ ತಾವೇ ತಿಳಿಯುತ್ತವೆಯೋ ಆ ಮಹತ್ವವಾದ ಜ್ಞಾನಕ್ಕೆ ನಮ್ಮ ಶಾಸ್ತ್ರದಲ್ಲಿ “ಬ್ರಹ್ಮವಿದ್ಯೆ”  ಎಂದು  ವಿವರಿಸಲಾಗಿದೆ – ಒಂದೇ ಒಂದರ ಜ್ಞಾನ, ಸಕಲ ಜ್ಞಾನವನ್ನೂ ತಿಳಿಯುವುದರೆಡೆಗೆ ಕೊಂಡೊಯ್ಯುತ್ತದೆ!
\end{inspiration}
\newpage

\slcol{\Index{ಮನುಷ್ಯಾಣಾಂ ಸಹಸ್ರೇಷು} ಕಶ್ಚಿದ್ಯತತಿ ಸಿದ್ಧಯೇ ।\\
ಯತತಾಮಪಿ ಸಿದ್ಧಾನಾಂ ಕಶ್ಚಿನ್ಮಾಂ ವೇತ್ತಿ ತತ್ತ್ವತಃ ॥ ೩ ॥}
\cquote{ಸಾವಿರ ಮನುಷ್ಯರಲ್ಲಿ ಒಬ್ಬನು ಸಿದ್ಧಿಗಾಗಿ ಯತ್ನಿಸುತ್ತಾನೆ. ಹಾಗೆ ಪ್ರಯತ್ನಿಸುತ್ತಾ ಸಿದ್ಧಿ ಪಡೆದ ಸಾವಿರ ಮಂದಿಯಲ್ಲಿ ಯಾವನೋ ಒಬ್ಬನು ಮಾತ್ರ ನನ್ನನ್ನು ಸರಿಯಾಗಿ ತಿಳಿಯುತ್ತಾನೆ.}
\slcol{\Index{ಭೂಮಿರಾಪೋಽನಲೋ} ವಾಯುಃ ಖಂ ಮನೋ ಬುದ್ಧಿರೇವ ಚ ।\\
ಅಹಂಕಾರ ಇತೀಯಂ ಮೇ ಭಿನ್ನಾ ಪ್ರಕೃತಿರಷ್ಟಧಾ ॥ ೪ ॥}
\cquote{ಭೂಮಿ,ನೀರು, ಬೆಂಕಿ, ಗಾಳಿ, ಆಕಾಶ, ಮನಸ್ಸು, ಬುದ್ಧಿ, ಅಹಂಕಾರ ಹೀಗೆ ನನ್ನ ಅಧೀನವಾದ ಜಡ ಪ್ರಕೃತಿಯು ಎಂಟು ಬಗೆಯಾಗಿರುವುದು.}
\slcol{\Index{ಅಪರೇಯಮಿತಸ್ತ್ವನ್ಯಾಂ} ಪ್ರಕೃತಿಂ ವಿದ್ಧಿ ಮೇ ಪರಾಮ್ ।\\
ಜೀವಭೂತಾಂ ಮಹಾಬಾಹೋ ಯಯೇದಂ ಧಾರ್ಯತೇ ಜಗತ್ ॥ ೫ ॥ }
\cquote{ಮಹಾಬಾಹೋ, ಇದು ಪ್ರಕೃತಿಯ ಕೀಳುಭಾಗ. ಇದಕ್ಕಿಂತ ಮೇಲೆ ಇನ್ನೊಂದು ಪ್ರಕೃತಿ ಇದೆ. ಅದು ಚೇತನ ಪ್ರಕೃತಿಯಾಗಿದೆ, ಈ ಜಗತ್ತನ್ನು ಹಿಡಿದು ನಿಲ್ಲಿಸಿಕೊಂಡಿದೆ.}
\slcol{\Index{ಏತದ್ಯೋನೀನಿ ಭೂತಾನಿ} ಸರ್ವಾಣೀತ್ಯುಪಧಾರಯ ।\\
ಅಹಂ ಕೃತ್ಸ್ನಸ್ಯ ಜಗತಃ ಪ್ರಭವಃ ಪ್ರಲಯಸ್ತಥಾ ॥ ೬ ॥}
\cquote{ಜಗತ್ತಿನಲ್ಲಿರುವುದೆಲ್ಲ ಈ ಎರಡರಿಂದ ಹೊರ ಹೊಮ್ಮಿದವುಗಳೆಂದು ತಿಳಿ. ನಾನು ಎಲ್ಲ ಜಗತ್ತಿನ ಸೃಷ್ಟಿಗೂ, ನಾಶಕ್ಕೂ ಕಾರಣನು.}
\slcol{\Index{ಮತ್ತಃ ಪರತರಂ} ನಾನ್ಯತ್ಕಿಂಚಿದಸ್ತಿ ಧನಂಜಯ ।\\
ಮಯಿ ಸರ್ವಮಿದಂ ಪ್ರೋತಂ ಸೂತ್ರೇ ಮಣಿಗಣಾ ಇವ ॥ ೭ ॥}
\cquote{ಧನಂಜಯ, ನನಗಿಂತ ಉತ್ತಮವಾದದ್ದು ಏನೂ ಇಲ್ಲ. ನೂಲಿನಲ್ಲಿ ಮಣಿಗಳು ಪೋಣಿಸಿರುವಂತೆ ನನ್ನಲ್ಲೇ  ಈ ಜಗತ್ತೆಲ್ಲಾ  ಹೆಣೆದುಕೊಂಡಿದೆ.}
\slcol{\Index{ರಸೋಽಹಮಪ್ಸು} ಕೌಂತೇಯ ಪ್ರಭಾಸ್ಮಿ ಶಶಿಸೂರ್ಯಯೋಃ ।\\
ಪ್ರಣವಃ ಸರ್ವವೇದೇಷು ಶಬ್ದಃ ಖೇ ಪೌರುಷಂ ನೃಷು ॥ ೮ ॥}
\cquote{ಕುಂತೀಪುತ್ರ, ನೀರಿನಲ್ಲಿ ರಸ ನಾನು, ಚಂದ್ರ ಸೂರ್ಯರಲ್ಲಿ ಪ್ರಭೆ, ವೇದಗಳಲೆಲ್ಲಾ ಓಂಕಾರ, ಆಕಾಶದಲ್ಲಿ ಧ್ವನಿ, ಪುರುಷರಲ್ಲಿ ಪೌರುಷ.}

\newpage
\begin{mananam}{\mananamTitle{ಮನನ ಶ್ಲೋಕ - ೩}}
\nfntmtxt ನಮ್ಮದೇ ಸಮಾಜದಲ್ಲಿ ಮತ್ತು ಸುತ್ತಮುತ್ತಲಿನ ಅನೇಕ ಜನರು ಭೌತಿಕವಾದ ಗುರಿಗಳನ್ನೇ ಕೇಂದ್ರವಾಗಿಟ್ಟುಕೊಂಡು ಅದರ ಬೆನ್ನು ಹತ್ತುವುದನ್ನು ನಾನು ಗಮನಿಸಿದ್ದೇನೆಯೇ? ಇತರರು ಪ್ರಾಥಮಿಕವಾಗಿ ತಮ್ಮದೇ ಕ್ಷೇತ್ರಗಳಲ್ಲಿ ವೈಯಕ್ತಿಕ ಯಶಸ್ಸನ್ನು ಸಾಧಿಸುವದರಲ್ಲಿಯೇ ಪ್ರಮುಖವಾಗಿ ನಿರತರಾಗಿದ್ದಾರೆ ಎಂಬುದನ್ನು ನಾನು ನೋಡಬಹುದೇ? ಕೇವಲ ಕೆಲವರು ಮಾತ್ರ ಏಕೆ ನಿಜವಾದ ಆತ್ಮಾಭಿವೃದ್ಧಿಯನ್ನು ಮತ್ತು ಆಧ್ಯಾತ್ಮಿಕ ಬೆಳವಣಿಗೆಯನ್ನು ಬೆನ್ನು ಹತ್ತಿ ಹೋಗುತ್ತಾರೆ? ತಮ್ಮ ನಿಜ ಸ್ವರೂಪವನ್ನು ಅನ್ವೇಷಿಸಲು ಬಯಸುವವರು ಅಥವಾ ದೈವೀ ಸಂಪರ್ಕ ಹೊಂದಲು ಪ್ರಯತ್ನ ಮಾಡುತ್ತಿರುವ ಕೆಲವೇ ಕೆಲವರಲ್ಲಿಯೂ ಸಹ, ಕೆಲವೇ ಕೆಲವರು ಮಾತ್ರ ಏಕೆ ಅಂತ್ಯದವರೆಗೆ ತಮ್ಮ ಅವಿರತ ಪ್ರಯತ್ನ ನಿಭಾಯಿಸುತ್ತಾರೆ? ನನ್ನ ಸ್ವಂತ ಜೀವನದಲ್ಲಿ, ಶಾಶ್ವತ ಸಂತೋಷಕ್ಕಾಗಿ ಉನ್ನತವಾದ ಮಾರ್ಗವನ್ನು ಆರಿಸುವುದನ್ನು ತಡೆಯುವ ಅಡೆತಡೆಗಳು ಯಾವುವು?
\end{mananam}
\sReflect
\begin{inspiration}{\mananamTitle{ಸ್ಪೂರ್ತಿ}}
\nfntmtxt \footnotesize ಉಪನಿಷತ್ ಒಂದರಲ್ಲಿ, ಪ್ರತಿಯೊಬ್ಬ ಮಾನವನಿಗೂ ತನ್ನ ಜೀವನದ ದಾರಿಯನ್ನು ಕಂಡುಕೊಳ್ಳುವ ಆಯ್ಕೆ ಇರುತ್ತದೆ ಎಂದು ಬೋಧಿಸಿದ್ದಾರೆ. ಕೆಲವರು ಇಂದ್ರಿಯಗಳನ್ನು ಮೆಚ್ಚಿಸುವ ಮತ್ತು ಅಹಂಕಾರವನ್ನು ಪೋಷಿಸುವ ಮಾರ್ಗವನ್ನು ಆರಿಸಿಕೊಳ್ಳುತ್ತಾರೆ -ಈ ಮಾರ್ಗವು ‘ಪ್ರೇಯಸ್’ ಅನಿಸಿಕೊಳ್ಳುತ್ತದೆ; ಇತರರು, ಉನ್ನತವಾದ ಉತ್ತಮ ಹಾದಿಗೆ ಬದ್ಧರಾಗಿರುತ್ತಾರೆ, ಇದು ಅಂತಿಮವಾಗಿ ಸಂತೋಷಕ್ಕೆ ಕಾರಣವಾಗುತ್ತದೆ -ಈ ಮಾರ್ಗವು “ಶ್ರೇಯಸ್” ಅನಿಸಿಕೊಳ್ಳುತ್ತದೆ; ಅನೇಕರಿಗೆ ಈ ಮಾರ್ಗಗಳ ಬಗ್ಗೆ ತಿಳಿದಿದ್ದರೂ ಸಹ, ಕೆಲವರು ಮಾತ್ರ ಎರಡನೆಯ ಮಾರ್ಗದಲ್ಲಿ ನಡೆಯುವ ಆಯ್ಕೆ ಮಾಡಿಕೊಳ್ಳುತ್ತಾರೆ; ಇವರ ಮಧ್ಯೆಯೂ ಕೂಡ ಕೆಲವರು ಮಾತ್ರ ಛಲದಿಂದ ಬಿಡದೆಲೇ ಮುಂದುವರೆಯುತ್ತಾರೆ; ಇಂತಹ ಸಂದರ್ಭದಲ್ಲಿಯೇ ಧರ್ಮಗ್ರಂಥಗಳು, ಗುರುಗಳ ಮಾರ್ಗದರ್ಶನ ಮತ್ತು ಸರಿಯಾದ ಸಹಚರರ ಬೆಂಬಲವು ಅನ್ವೇಷಕನಿಗೆ (ಸಾಧಕ) ಅಗತ್ಯವಾಗುವುದು; ಅವರೇ ನಿಮ್ಮ ಅತ್ಯಂತ ನಿಷ್ಠಾವಂತ    ಪ್ರೋತ್ಸಾಹಕರು!
\end{inspiration}
\newpage

\begin{mananam}{\mananamTitle{ಮನನ ಶ್ಲೋಕ - ೪, ೫}}
\nfntmtxt ನಾವು ಅನುಭವಿಸುವ ಪ್ರಕೃತಿಯ ವಿವಿಧ ಶಕ್ತಿಗಳು ಯಾವುವು?  ವೈಜ್ಞಾನಿಕ ದೃಷ್ಟಿಕೋನದಿಂದ, ನಾವು ಅವುಗಳನ್ನು ವಸ್ತು, ಶಕ್ತಿ ಮತ್ತು ಮನಸ್ಸು ಎಂದೂ ಸಹ ವರ್ಗೀಕರಿಸುತ್ತೇವೆ. ಪ್ರಕೃತಿಯಲ್ಲಿರುವ ಈ ಶಕ್ತಿಗಳು ಸ್ವತಂತ್ರವಾಗಿವೆಯೇ ಅಥವಾ, ಅವು ಪರಸ್ಪರ ಅವಲಂಬಿತವಾಗಿವೆಯೇ? ಆದಾಗ್ಯೂ, ಇವುಗಳನ್ನು ಪೋಷಿಸುವಂತಹ ಹಾಗೂ ವಿಭಿನ್ನವಾಗಿರುವ ಬೇರೆಯೇ ಒಂದು ಶಕ್ತಿ ಇದೆಯೇ? ವಿಜ್ಞಾನ ಅಥವಾ ಧರ್ಮದ ಪರಿಕಲ್ಪನೆಗಳನ್ನು ಅವಲಂಬಿಸದೇ, ಪ್ರಕೃತಿಯಲ್ಲಿ ಈ ವಿಭಿನ್ನ ಶಕ್ತಿಗಳ ಕಾರ್ಯಗಳನ್ನು ನಾನು ನೇರವಾಗಿ ಗಮನಿಸಬಹುದೇ?
\end{mananam}
\sReflect
\begin{inspiration}{\mananamTitle{ಸ್ಪೂರ್ತಿ}}
\nfntmtxt ಪ್ರಕೃತಿಯಲ್ಲಿ “ಪರ” ಮತ್ತು “ಅಪರ” ಎಂಬ ಎರಡು ಶಕ್ತಿಗಳು ಕಾರ್ಯಯುಕ್ತವಾಗಿವೆ ಎಂದು ಗೀತೆಯು ಸೂಚಿಸುತ್ತದೆ. ಮನಸ್ಸು ಮತ್ತು ವಸ್ತುವಿನ ಕಾರ್ಯಗಳು ಎಲ್ಲರಿಗೂ ಸ್ಪಷ್ಟವಾಗಿದ್ದರೂ ಸಹ, ಎಲ್ಲಕ್ಕೂ ಆಧಾರವಾಗಿರುವ ಮತ್ತು ಎಲ್ಲವನ್ನೂ ಒಟ್ಟಿಗೆ ಹಿಡಿದಿಟ್ಟಿರುವ ಆ  ‘ದಿವ್ಯ ಚೈತನ್ಯ’ ದ  ಪಾತ್ರವು ಮಾನವರು ಬಿಚ್ಚಿಡಲು ಉದ್ದೇಶಿಸಿರುವ ಆಳವಾದ ರಹಸ್ಯವಾಗಿಯೇ ಉಳಿಯುತ್ತದೆ. ಎಲ್ಲಿಯವರೆಗೆ ನಾವು ವಸ್ತುವಿನೊಂದಿಗೆ ಗುರುತಿಸಲ್ಪಟ್ಟ ಜೀವನದೊಂದಿಗೆ ತೃಪ್ತರಾಗುತ್ತೇವೋ ಅಥವಾ ಹೆಚ್ಚೆoದರೆ, ಮನಸ್ಸಿನಲ್ಲಿಯೇ ತೃಪ್ತಿ ಕಾಣುತ್ತೇವೆಯೋ ಅಲ್ಲಿಯವರೆಗೂ, ನಮ್ಮ ನಿಜವಾದ ಅಸ್ತಿತ್ವವು ನಮ್ಮಿಂದ ಮರೆಮಾಚಲ್ಪಡುತ್ತದೆ. ಒಬ್ಬ ‘ಜಿಜ್ಞಾಸುವು’ ತನ್ನ ಆಳವಾದ ಅಸ್ತಿತ್ವವನ್ನು ಅರಿಯಲು ಪ್ರಾಮಾಣಿಕವಾಗಿ ಅಪೇಕ್ಷಿಸುತ್ತಾನೆ; ಆದರೆ,  ಒಬ್ಬ ‘ಮುಮುಕ್ಷು’ವು ಯಾವುದೇ ಅಡೆತಡೆಯನ್ನೂ ಸಹ ಲೆಕ್ಕಿಸದೆ ದೃಢವಾಗಿ  ಅನುಸರಿಸಿ ಹೋಗಲು ನಿಶ್ಚಯಿಸುತ್ತಾನೆ.
\end{inspiration}
\newpage

\begin{mananam}{\mananamTitle{ಮನನ ಶ್ಲೋಕ - ೭}}
\nfntmtxt ಪ್ರಕೃತಿಯ ಇತರ ಎಲ್ಲಾ ಶಕ್ತಿಗಳನ್ನೂ ಆಳುವ ಸರ್ವೋಚ್ಚ ಶಕ್ತಿಯನ್ನು ನಾನು ಏನೆಂದು ಅರ್ಥಮಾಡಿಕೊಂಡಿದ್ದೇನೆ?  ಎಲ್ಲವನ್ನೂ ಪೋಷಿಸುವ, ಸಮಸ್ತ ಸೃಷ್ಟಿಯ ಅಡಿಪಾಯವಾಗಿ ಕಾರ್ಯನಿರ್ವಹಿಸುವ ಈ ಅಗತ್ಯ ಶಕ್ತಿಯನ್ನು ನಾನು ಹೇಗೆ ಗ್ರಹಿಸುವುದು? ಇದರ ಬಗ್ಗೆ ಸನಾತನ ಗ್ರಂಥಗಳು  ಯಾವ ರೀತಿಯ ಒಳನೋಟಗಳನ್ನು ಒದಗಿಸುತ್ತವೆ? ನಾನು ಎಲ್ಲಿ ಮತ್ತು ಹೇಗೆ ಇದನ್ನು ನೋಡಲಿ, ತಿಳಿಯಲಿ ಅಥವಾ ಗ್ರಹಿಸಲಿ? ನನ್ನ ಸ್ವಂತ ದೇಹ ಮತ್ತು ಮನಸ್ಸಿನೊಳಗೆ ನಾನು ಈ ಪ್ರಮುಖವಾದ ಶಕ್ತಿಯ ಕಾರ್ಯಗಳನ್ನು ಹೇಗೆ ಗ್ರಹಿಸಲಿ? ; ಯಾವ ಆಧ್ಯಾತ್ಮಿಕ ಅಭ್ಯಾಸಗಳು ಇದರ ಅನುಭವವನ್ನು ನೇರವಾಗಿ ಪಡೆಯಲು ನನಗೆ ಸಹಾಯಕವಾಗಬಲ್ಲದು?
\end{mananam}
\sReflect
\begin{inspiration}{\mananamTitle{ಸ್ಪೂರ್ತಿ}}
\nfntmtxt ಕಂಠಹಾರದಲ್ಲಿರುವ ಮುತ್ತುಗಳನ್ನು ಒಂದು ದಾರವು ಹಿಡಿದಿಟ್ಟಿರುತ್ತದೆ, ಆದರೆ ಗಮನ ಸೆಳೆಯುವುದು ಆ ಮುತ್ತುಗಳೇ ಹೊರತು, ದಾರವಲ್ಲ; ಅದೇ ರೀತಿ ಭಗವಂತ ಒಬ್ಬ ಕಣ್ಣಿಗೆ ಕಾಣದ ಶಕ್ತಿ, ಬ್ರಹ್ಮಾಂಡದ ಎಲ್ಲಾ ಶಕ್ತಿಗಳ ಕೆಲಸದ ಮೂಲ;  ನೇರವಾಗಿ ಅನುಭವಕ್ಕೆ ಬರದಿದ್ದರೂ, ಹೇಗೆ ಗುರುತ್ವಾಕರ್ಷಣ ಶಕ್ತಿಯು ಸದಾ ಕೆಲಸ ಮಾಡುತ್ತಿರುವುದೋ ಹಾಗೆಯೇ, ಕಣ್ಣಿಗೆ ಕಾಣದಿದ್ದರೂ, ಆಧ್ಯಾತ್ಮದ ಅರಿವು ಇರುವವರಿಗೆ ಈ ದೈವಿಕ ಶಕ್ತಿಯು ಎಲ್ಲದರ ಮೂಲತತ್ವ.
\end{inspiration}
\newpage

\slcol{\Index{ಪುಣ್ಯೋ ಗಂಧಃ} ಪೃಥ್ವಿವ್ಯಾಂ ಚ ತೇಜಶ್ಚಾಸ್ಮಿ ವಿಭಾವಸೌ ।\\
ಜೀವನಂ ಸರ್ವಭೂತೇಷು ತಪಶ್ಚಾಸ್ಮಿ ತಪಸ್ವಿಷು ॥ ೯ ॥ }
\cquote{ಭೂಮಿಯಲ್ಲಿ  ಕಂಪು ನಾನು, ಬೆಂಕಿಯಲ್ಲಿ ಬೆಳಕು ನಾನು, ಎಲ್ಲ ಪ್ರಾಣಿಗಳಲ್ಲಿ ಬದುಕು ನಾನು, ತಪಸ್ವಿಗಳಲ್ಲಿ ತಪಸ್ಸು ನಾನು. }
\slcol{\Index{ಬೀಜಂ ಮಾಂ} ಸರ್ವಭೂತಾನಾಂ ವಿದ್ಧಿ ಪಾರ್ಥ ಸನಾತನಮ್ ।\\
ಬುದ್ಧಿರ್ಬುದ್ಧಿಮತಾಮಸ್ಮಿ ತೇಜಸ್ತೇಜಸ್ವಿನಾಮಹಮ್ ॥ ೧೦ ॥ }
\cquote{ಪಾರ್ಥ, ನನ್ನನ್ನು ಎಲ್ಲ ಚರಾಚರಗಳ ಸನಾತನವಾದ ಬೀಜಶಕ್ತಿ ಎಂದು ತಿಳಿ. ನಾನು ಬುದ್ಧಿವಂತರಲ್ಲಿ ವಿವೇಕಶಕ್ತಿಯಾಗಿದ್ದೇನೆ. ಆತ್ಮ ಶಕ್ತಿಯುಳ್ಳವರಲ್ಲಿ ಆತ್ಮಶಕ್ತಿ ಆಗಿದ್ದೇನೆ.}
\slcol{\Index{ಬಲಂ ಬಲವತಾಂ} ಚಾಹಂ ಕಾಮರಾಗವಿವರ್ಜಿತಮ್ ।\\
ಧರ್ಮಾವಿರುದ್ಧೋ ಭೂತೇಷು ಕಾಮೋಽಸ್ಮಿ ಭರತರ್ಷಭ ॥ ೧೧ ॥}
\cquote{ಭರತಕುಮಾರ, ಬಲವಂತರಲ್ಲಿ ಆಸೆ ಆಸಕ್ತಿಗಳಿಗೆ ಬಲಿಯಾಗದ ಬಲ ನಾನು. ಪ್ರಾಣಿಗಳಲ್ಲಿ ಧರ್ಮ ಸಮ್ಮತವಾದ ಬಯಕೆ ನಾನಾಗಿದ್ದೇನೆ.}
\slcol{\Index{ಯೇ ಚೈವ ಸಾತ್ತ್ವಿಕಾ} ಭಾವಾ ರಾಜಸಾಸ್ತಾಮಸಾಶ್ಚ ಯೇ ।\\
ಮತ್ತ ಏವೇತಿ ತಾನ್ವಿದ್ಧಿ ನ ತ್ವಹಂ ತೇಷು ತೇ ಮಯಿ ॥ ೧೨ ।}
\cquote{ಇಷ್ಟೇ ಎಂದಲ್ಲ. ಸಾತ್ವಿಕವಾದ, ರಾಜಸವಾದ, ತಾಮಸವಾದ ಏನೆಲ್ಲ ವಸ್ತುಸ್ವಭಾವಗಳಿವೆಯೋ ಅವೆಲ್ಲ ನನ್ನಿಂದಲೇ ಹುಟ್ಟಿದವುಗಳೆಂದು ತಿಳಿ. ನಾನು ಅವುಗಳ ಆಸರೆಯಲಿಲ್ಲ, ಅವು ನನ್ನ  ಆಸರೆಯಲ್ಲಿವೆ.}
\slcol{\Index{ತ್ರಿಭಿರ್ಗುಣಮಯೈರ್ಭಾವೈರೇಭಿಃ} ಸರ್ವಮಿದಂ ಜಗತ್ ।\\
ಮೋಹಿತಂ ನಾಭಿಜಾನಾತಿ ಮಾಮೇಭ್ಯಃ ಪರಮವ್ಯಯಮ್ ॥ ೧೩ ॥}
\cquote{ಈ ಮೂರು ಬಗೆಯ ಗುಣಗಳಿಂದ ರೂಪುಗೊಂಡ ವಸ್ತುಗಳಿಂದ ಮೋಹಗೊಂಡ ಈ ಎಲ್ಲ ಜಗತ್ತು, ಈ ತ್ರಿಗುಣಗಳಿಂದ ಆಚೆಗಿರುವ ನನ್ನ ನಿತ್ಯ ಸ್ವರೂಪವನ್ನು ಅರಿಯಲಾರದಾಗಿದೆ.}


\newpage
\begin{mananam}{\mananamTitle{ಮನನ ಶ್ಲೋಕ - ೧೦}}
\nfntmtxt ಸಣ್ಣ ಬೀಜವು ಹೆಮ್ಮರವಾಗಿ ಬೆಳೆಯಲು ಅನುವು ಮಾಡಿಕೊಡುವುದಂತಹದ್ದು ಆ ಬೀಜದಲ್ಲಿ ಏನಿದೆ? ಆಲೋಚನೆ ಮತ್ತು ಅಭಿವ್ಯಕ್ತಿಯನ್ನು ಉಂಟುಮಾಡುವ ಸಾಮರ್ಥ್ಯವನ್ನು ಮೆದುಳಿಗೆ  ನೀಡುವುದು ಯಾವುದು?  ಸೂರ್ಯನ ಬೆಳಕು ಮತ್ತು ನಕ್ಷತ್ರಗಳ ಮಿನುಗುವಿಕೆಯನ್ನು ಯಾವುದು ಬೆಳಗಿಸುತ್ತದೆ? ಪ್ರಕೃತಿಯ ಅತ್ಯಂತ ಭವ್ಯವಾದ ಅದ್ಭುತಗಳಲ್ಲಿ ಸೌಂದರ್ಯವನ್ನು ತುಂಬಿರುವವರಾರು? ನಾನು, ಒಂದು ಮಗುವಿನಲ್ಲಿರುವ ಕುತೂಹಲ ಮತ್ತು ವಿಸ್ಮಯಗಳನ್ನು ಕಳೆದುಕೊಂಡಿದ್ದೇನೆಯೆ? ವಿಸ್ಮಯದಿಂದ ಪ್ರಕೃತಿಯನ್ನು ನೋಡುವುದರಲ್ಲಿ ಮೌಲ್ಯವಿದೆಯೇ? ಇದರತ್ತ ತೋರುವ ಮೆಚ್ಚುಗೆಯು, ಅದರ ಹಿಂದೆ ಅಡಗಿಕೊಂಡಿರುವ ದೈವತ್ವವನ್ನು ಕಂಡುಕೊಳ್ಳುವುದರತ್ತ ನನಗೆ ದಾರಿ ತೋರಿಸುತ್ತದೆಯೆ?
\end{mananam}
\sReflect
\begin{inspiration}{\mananamTitle{ಸ್ಪೂರ್ತಿ}}
\nfntmtxt ಸೌಂದರ್ಯ ಮತ್ತು ಜಟಿಲತೆಯಿಂದ ಕೂಡಿದ ಈ ಸೃಷ್ಟಿಯು ಕೇವಲ ಒಂದು ಆಕಸ್ಮಿಕವಾಗಿರಲು ಸಾಧ್ಯವಿಲ್ಲ. ಗ್ರಹಗಳು ಮತ್ತು ನಕ್ಷತ್ರಗಳನ್ನು ಅವುಗಳ ಸ್ಥಾನದಲ್ಲಿ ಹಿಡಿದಿಟ್ಟಿರುವ ಶಕ್ತಿ, ನಮ್ಮ ಜೀರ್ಣಕ್ರಿಯೆ ಮತ್ತು ಉಸಿರಾಟವನ್ನು ನಿರ್ವಹಿಸುವ ಈ ಶಕ್ತಿ, ಇವೆಲ್ಲವೂ ಸಕಲ ಅಸ್ತಿತ್ವದಲ್ಲಿಯೇ ಅತ್ಯಂತ ಉದಾತ್ತವಾದ  ‘ಇರುವಿಕೆ’ಯೇ ಆಗಿರುತ್ತದೆ. ಈ ‘ಇರುವಿಕೆ’ಯು ತನ್ನ ಸೇವೆಗೆ ಬೆಲೆಯನ್ನು ಕೇಳುವುದಾಗಿದ್ದಲ್ಲಿ, ನಾನು ಶಾಶ್ವತವಾಗಿ ಅದರ ಅಡಿಯಾಳಾಗಿರಬೇಕಾಗಿರುತ್ತಿತ್ತು! ಆದಾಗ್ಯೂ, ಪ್ರತಿಯಾಗಿ ನಮ್ಮಿಂದ ಏನನ್ನೂ ನಿರೀಕ್ಷಿಸದೆ ದೇವರು ನಮಗೆ ಸೇವೆ ಮಾಡುತ್ತಿದ್ದಾನೆ. ಈ ರೀತಿಯ ನಿಬಂಧನ ರಹಿತವಾದ ಪ್ರೀತಿಗೆ ನಾವು ಕನಿಷ್ಠಪಕ್ಷ, ನಮ್ಮ ಕೃತಜ್ಞತೆಯನ್ನಾದರೂ ವ್ಯಕ್ತಪಡಿಸಬೇಕಲ್ಲವೇ?
\end{inspiration}
\newpage

\slcol{\Index{ದೈವೀ ಹ್ಯೇಷಾ ಗುಣಮಯೀ} ಮಮ ಮಾಯಾ ದುರತ್ಯಯಾ ।\\
ಮಾಮೇವ ಯೇ ಪ್ರಪದ್ಯಂತೇ ಮಾಯಾಮೇತಾಂ ತರಂತಿ ತೇ ॥ ೧೪ ॥}
\cquote{ದಿವ್ಯವೂ, ಗುಣಮಯವೂ ಆದ ನನ್ನ ಈ ಮಾಯಾಶಕ್ತಿಯನ್ನು ದಾಟುವುದು ಕಷ್ಟ.ಯಾರು ನನಗೆ ಶರಣು ಬರುತ್ತಾರೋ ಅವರು ಈ ಮಾಯೆಯನ್ನು ದಾಟಬಲ್ಲರು.}
\slcol{\Index{ನ ಮಾಂ ದುಷ್ಕೃತಿನೋ} ಮೂಢಾಃ ಪ್ರಪದ್ಯಂತೇ ನರಾಧಮಾಃ ।\\
ಮಾಯಯಾಪಹೃತಜ್ಞಾನಾ ಆಸುರಂ ಭಾವಮಾಶ್ರಿತಾಃ ॥ ೧೫ ॥}
\cquote{ಮಾಯೆಯ ಮೋಹಕ್ಕೆ ಒಳಗಾಗಿ ಅರಿವುಗೇಡಿಗಳಾದ, ದುಷ್ಟ ಸ್ವಭಾವದ ದುಷ್ಕರ್ಮಿಗಳು, ಮೂಢರಾದ ಹೀನ ಮಾನವರು ನನ್ನನ್ನು ಶರಣು ಹೊಂದುವುದಿಲ್ಲ.}
\slcol{\Index{ಚತುರ್ವಿಧಾ ಭಜಂತೇ} ಮಾಂ ಜನಾಃ ಸುಕೃತಿನೋಽರ್ಜುನ ।\\
ಆರ್ತೋ ಜಿಜ್ಞಾಸುರರ್ಥಾರ್ಥೀ ಜ್ಞಾನೀ ಚ ಭರತರ್ಷಭ ॥ ೧೬ ॥}
\cquote{ಅರ್ಜುನ, ದುಃಖಕ್ಕೆ ಒಳಗಾದವನು, ತಿಳಿಯ ಬಯಸುವವನು, ಸುಖ ಸಾಧನಗಳನ್ನು ಬಯಸುವವನು, ಆತ್ಮ ಸ್ವರೂಪವನ್ನರಿತವನು, ಈ ನಾಲ್ಕು ಬಗೆಯ ಪುಣ್ಯಶೀಲರಾದ ಜನರು ನನ್ನನ್ನು ಭಜಿಸುತ್ತಾರೆ.}
\slcol{\Index{ತೇಷಾಂ ಜ್ಞಾನೀ} ನಿತ್ಯಯುಕ್ತ ಏಕಭಕ್ತಿರ್ವಿಶಿಷ್ಯತೇ ।\\
ಪ್ರಿಯೋ ಹಿ ಜ್ಞಾನಿನೋಽತ್ಯರ್ಥಮಹಂ ಸ ಚ ಮಮ ಪ್ರಿಯಃ ॥ ೧೭ ॥}
\cquote{ಅವರಲ್ಲಿ, ಯಾವಾಗಲೂ ನನ್ನನ್ನೇ ಹೊಂದಿಕೊಂಡು ನನ್ನನ್ನೇ ಪ್ರೀತಿಸುವ ಜ್ಞಾನಿಯು ಹೆಚ್ಚಿನವನು. ಏಕೆಂದರೆ ನಾನು ಜ್ಞಾನಿಗೆ ಬಹಳ ಪ್ರೀತಿಯ ಸೊತ್ತು. ಅವನೂ ನನಗೆ ಅತ್ಯಂತ ಪ್ರಿಯ.}
\slcol{\Index{ಉದಾರಾಃ ಸರ್ವ ಏವೈತೇ} ಜ್ಞಾನೀ ತ್ವಾತ್ಮೈವ ಮೇ ಮತಮ್ ।
ಆಸ್ಥಿತಃ ಸ ಹಿ ಯುಕ್ತಾತ್ಮಾ ಮಾಮೇವಾನುತ್ತಮಾಂ ಗತಿಮ್ ॥ ೧೮ ॥}
\cquote{ಈ ಎಲ್ಲರೂ ಮೇಲಾದವರೇ. ಜ್ಞಾನಿಯಾದರೋ ನನ್ನ ಅಂತರಂಗದ ಸ್ವತ್ತು ಎಂದು ನನ್ನ ಅಭಿಪ್ರಾಯ. ಏಕೆಂದರೆ ಅವನು ಮನಸ್ಸನ್ನು ನನ್ನಲ್ಲಿ ನೆಲೆಗೊಳಿಸಿ ಹಿರಿಯ ಆಸರೆಯಾದ ನನ್ನನ್ನೇ ನಂಬಿರುವವನು.}

\newpage
\begin{mananam}{\mananamTitle{ಮನನ ಶ್ಲೋಕ - ೧೪}}
\nfntmtxt ಮೋಹದ(ಭ್ರಮೆಯ) ಗುಣಲಕ್ಷಣಗಳೇನು? ಹಿಂದೂ ಸನಾತನ ಗ್ರಂಥಗಳು   ತಿಳಿಸಿದಂತೆ ನಾವೆಲ್ಲ ಸಾಮೂಹಿಕ ಭ್ರಮೆಯಲ್ಲಿ  ಸಿಕ್ಕಿಹಾಕಿಕೊಂಡಿದ್ದೇವೆಯೇ? ಅಂತ್ಯದಲ್ಲಿ ವಿಷಾದವೇ ಪಡುವಂತಹ, ಭ್ರಮಾಧೀನ ಮನೋಸ್ಥಿತಿಯಲ್ಲಿ, ಯಾವ ಯಾವ ಕಾರ್ಯಗಳನ್ನು ನನ್ನ ಜೀವನದಲ್ಲಿ ಮಾಡಿದ್ದೇನೆ? ನನ್ನನ್ನು ಜಾಗೃತ ಸ್ಥಿತಿಗೆ ಮರಳಿಸುವುದಕ್ಕೆ ಯಾವುದು ನನಗೆ ಸಹಾಯಕವಾಗಬಹುದು? ಮೋಹಯುಕ್ತ ಮನಸ್ಸಿನಿಂದ ಹುಟ್ಟುವ  ಇನ್ನೂ ಯಾವ ಯಾವ ಅಭ್ಯಾಸಗಳನ್ನು ನಾನು ರೂಢಿಸಿಕೊಂಡಿದ್ದೇನೆ? ಅಂತಹ ಮೋಹ ಜಾಲದಲ್ಲಿ ಪುನಃ ಸಿಲುಕಿಕೊಳ್ಳುವುದನ್ನು  ನಾನು ಹೇಗೆ ತಡೆಯಬಹುದು? ಇಂತಹ ಸ್ಥಿತಿಯಿಂದ ಮುಕ್ತವಾಗಲು ನನಗೆ ಸಹಾಯ ಮಾಡುವ ಯಾವುದೇ ಬಾಹ್ಯ ಮಾರ್ಗದರ್ಶನ ಅಥವಾ ಬೆಂಬಲವಿದೆಯೇ? 
\end{mananam}
\sReflect
\begin{inspiration}{\mananamTitle{ಸ್ಪೂರ್ತಿ}}
\nfntmtxt ಮಾಯಾ ಪರಿಕಲ್ಪನೆಯು, ವಿಶ್ವದಾದ್ಯಂತ ತತ್ತ್ವಚಿಂತನೆಗಳ ಅತ್ಯಂತ ಆಳವಾದ ವಿವರಣೆಗಳಲ್ಲಿ ಒಂದಾಗಿದೆ, ಇದು ನಮ್ಮ ನೈಜ ಸ್ವರೂಪದಿಂದ  ದೈವತ್ವವನ್ನು ಬೇರ್ಪಡಿಸುವ   ಶಕ್ತಿಯನ್ನು ವಿವರಿಸುತ್ತದೆ. ಆಧುನಿಕ ಉದಾಹರಣೆಗಳಾದ ಚಲನಚಿತ್ರಗಳು ಮತ್ತು ವಾಸ್ತವಕ್ಕೆ ನಿಕಟವಾದ ಆಧುನಿಕ ತಂತ್ರಜ್ಞಾನಗಳು  (ವರ್ಚುವಲ್ ರಿಯಾಲಿಟಿ)  ಮಾಯಾ ಕಲ್ಪನೆಯನ್ನು ಗ್ರಹಿಸಲು ಸ್ಪಷ್ಟವಾದ ಮಾರ್ಗವನ್ನು ಒದಗಿಸುತ್ತವೆ. ಹೇಗೆ ಒಬ್ಬ ಮಾಂತ್ರಿಕನು ತನ್ನ ಕೈಚಳಕ ಮತ್ತು ಹತೋಟಿಯಿಂದ ಇಂದ್ರಜಾಲ ಪ್ರದರ್ಶನವನ್ನು ನಡೆಸುತ್ತಾನೆಯೋ ಹಾಗೆಯೇ, ಮಾಯೆ ಎಂಬ ಭ್ರಾಂತಿಯ ಶಕ್ತಿಯನ್ನು ನಿಯಂತ್ರಿಸುವ ದೈವೀಕ ಮಾಂತ್ರಿಕನೊಬ್ಬನಿದ್ದಾನೆ; ಈ ಮಾಂತ್ರಿಕನ ಉದ್ದೇಶವು ಕ್ರೂರ ಅಥವಾ ಕುಯುಕ್ತಿಯನ್ನು ಆಧರಿಸಿರುವುದಿಲ್ಲ; ಎಲ್ಲಿ ಭ್ರಮೆಗಳು ಮತ್ತು ದುಃಸ್ವಪ್ನಗಳು ಶಾಶ್ವತವಾಗಿ ಕೊನೆಗೊಳ್ಳುತ್ತವೆಯೋ ಅಂತಹ ವಾಸ್ತವದಲ್ಲಿ, ಆತನ ಆಶ್ರಯ ಬಯಸಿ ಬಂದವರನ್ನು ಎತ್ತಿಹಿಡಿಯುತ್ತಾನೆ
\end{inspiration}
\newpage

\begin{mananam}{\mananamTitle{ಮನನ ಶ್ಲೋಕ - ೧೫}}
\nfntmtxt ಉನ್ನತವಾದ ಸತ್ಯವನ್ನು ಅರಿಯಲು ನನಗೆ ಯಾವುದು ಪ್ರೇರಣೆಯಾಯಿತು?ದೇವರನ್ನು ಹುಡುಕಲು, ಜೀವನದ ಯಾವ ಅನುಭವಗಳು  ನನಗೆ ಮಾರ್ಗದರ್ಶನ ನೀಡಿವೆ? 
ನನ್ನನ್ನು ಆಧ್ಯಾತ್ಮಿಕ ಹಾದಿಯತ್ತ ಕರೆದೊಯ್ಯಲು ಕಾರಣವಾಗಿರುವ, ನಾನು ಹಿಂದೆ ಪಟ್ಟಂತಹ ಸಂಕಟ ಅಥವಾ ನಷ್ಟ, ಅಥವಾ ಪೂಜ್ಯರ ಅನುಗ್ರಹ ಅಥವಾ ಮಾರ್ಗದರ್ಶನದ ಕ್ಷಣಗಳನ್ನು ನಾನು ಕೃತಜ್ಞತೆಯಿಂದ ಕಾಣುತ್ತೇನೆಯೇ?
ಸತ್ಯವನ್ನು ಹುಡುಕುವುದನ್ನು ಅಥವಾ ದೈವವನ್ನು ಗೌರವಿಸುವುದನ್ನು ನಿರ್ಲಕ್ಷಿಸುವಂತಹ ಭ್ರಮೆಯ ಸ್ಥಿತಿಗೆ ಮರಳುವುದನ್ನು ತಪ್ಪಿಸಲು ನಾನು, ನನ್ನ ವಿವೇಕವನ್ನು ಹೇಗೆ ಕಾಪಾಡಿಕೊಳ್ಳಬೇಕು?
\end{mananam}
\sReflect
\begin{inspiration}{\mananamTitle{ಸ್ಪೂರ್ತಿ}}
\nfntmtxt ಈ ಪ್ರಪಂಚದಲ್ಲಿ ಅನೇಕ ಜನರು ಎಂದಿಗೂ ಆಧ್ಯಾತ್ಮಿಕ ಅಭ್ಯಾಸಗಳನ್ನು ಅನ್ವೇಷಿಸದೇ ಅಥವಾ, ತನ್ನದೇ ಉನ್ನತ ಸತ್ಯವೇನೆಂದು ಅರಸದೇ ಸಂತೃಪ್ತರಾಗಿರುತ್ತಾರೆ.ಈ ವ್ಯಕ್ತಿಗಳನ್ನು ‘ಪಾಮರರು’ ಎಂದು ಕರೆಯುತ್ತಾರೆ ಮತ್ತು ಇವರು ಕೇವಲ ಇಂದ್ರಿಯ ಕಾಮನೆ ಪ್ರಚೋದನೆಯಿಂದಲೇ ತಮ್ಮ ಜೀವನ ನಡೆಸುತ್ತಿರುತ್ತಾರೆ. ಇಂತಹವರಿಗೆ ಉನ್ನತವಾದ ಸತ್ಯದ ಅರಿವಿನ ಕೊರತೆ ಇರುತ್ತದೆ ಮತ್ತು ಅದರ ಅನ್ವೇಷಣೆಯ ಬಗ್ಗೆ ಒಲವೂ ಸಹ ಇರುವುದಿಲ್ಲ. ಅವರ ಗಮನವು ಸಂಪೂರ್ಣವಾಗಿ ಪ್ರಾಪಂಚಿಕ  ವಿಷಯಗಳಲ್ಲಿ ಸಂತೋಷ ಕಂಡುಕೊಳ್ಳುವುದರ ಕಡೆಗೆ ಇರುತ್ತದೆ. ಅದಾಗ್ಯೂ, ಅವರ ಹೃದಯಾಂತರಾಳದಲ್ಲಿ  ಮರಣದ ಹಾಗೂ ಮರಣಾನಂತರ  ಜೀವನದ ಬಗ್ಗೆ ಸುಪ್ತವಾದ ಭಯ ಮನೆ ಮಾಡಿರುತ್ತದೆ ಮತ್ತು  ಯಾವಾಗಲೂ ಕಾಡುತ್ತಿರುತ್ತದೆ.
\end{inspiration}
\newpage


\begin{mananam}{\mananamTitle{ಮನನ ಶ್ಲೋಕ - ೧೬}}
\nfntmtxt ದೇವರ ಅಥವಾ ಸತ್ಯದ ಕಡೆಗೆ ನನ್ನ ಪ್ರಯಾಣದಲ್ಲಿ ನಾನು ಎಲ್ಲಿದ್ದೇನೆ? ನನ್ನ ಮುಖ್ಯ ಪ್ರೇರಣೆ ಏನು? ನಾನು  ಸಂಕಟದಲ್ಲಿದ್ದಾಗ ಅಥವಾ ಅವಶ್ಯಕತೆ ಇದ್ದಾಗ ಮಾತ್ರ ಉನ್ನತ ಶಕ್ತಿಯನ್ನು ಪ್ರಾರ್ಥಿಸುತ್ತೇನೆಯೆ? ಲೌಕಿಕ ಯಶಸ್ಸು ಮತ್ತು ರಕ್ಷಣೆಗಾಗಿ ನಾನು ಶ್ರೇಷ್ಠರಿಂದ ಆಶೀರ್ವಾದ ಕೋರುತ್ತೇನೆಯೇ? ನನಗೆ ದೈವೀಕ ಶಕ್ತಿಯನ್ನು ಅರಿಯುವ ಅಥವಾ, ನನ್ನ ಕಟ್ಟ ಕಡೆಯ ಅಸ್ತಿತ್ವದ ಮೂಲ ತತ್ವವನ್ನು ತಿಳಿಯುವ ಬಲವಾದ ಆಸೆ ಇದೆಯೇ?  ಉನ್ನತ ಜ್ಞಾನವನ್ನು ತಿಳಿಸಿಕೊಡುವವರ ಬಗ್ಗೆ ನನಗೆ ಗೌರವ   ಅಥವಾ ಭಕ್ತಿಯ ಭಾವನೆ ಇದೆಯೇ? ದೇವರ ಬಗೆಗಿನ ನನ್ನ ಭಕ್ತಿಯನ್ನು ಹೇಗೆ ಶುದ್ಧೀಕರಿಸಿಕೊಳ್ಳಲಿ ಮತ್ತು ಪರಮ ಸತ್ಯದ  ಬಗ್ಗೆ ನನ್ನ ತಿಳುವಳಿಕೆಯನ್ನು ಹೇಗೆ ಗಾಢವಾಗಿಸಲಿ?
\end{mananam}
\sReflect
\begin{inspiration}{\mananamTitle{ಸ್ಪೂರ್ತಿ}}
\nfntmtxt ದೇವರನ್ನು ಹುಡುಕುವವರಲ್ಲಿ ಸಹ, ಪ್ರೇರಣೆಗಳು ಬಹಳ ಭಿನ್ನವಾಗಿರುತ್ತವೆ. ಆರಂಭದಲ್ಲಿ, ಅನೇಕ ಅನ್ವೇಷಕರು ಕೆಲವು ರೀತಿಯ ದುಃಖಗಳನ್ನು ನಿವಾರಿಸುವ ಮತ್ತು ಲೌಕಿಕ ಸಂತೋಷವನ್ನು ಪಡೆಯುವ ಬಯಕೆಯಿಂದ ಪ್ರೇರೇಪಿಸಲ್ಪಡುತ್ತಾರೆ. ಕೆಲವರು, ಇತರರ ಮೇಲೆ ಸ್ಪರ್ಧಾತ್ಮಕ ರೀತಿಯಿಂದ ಮುಂದಿರಲೆಂದೇ  ಆಧ್ಯಾತ್ಮಿಕ ಅಭ್ಯಾಸಗಳನ್ನು ಕೈಗೊಳ್ಳುತ್ತಾರೆ. ಅಂತಿಮವಾಗಿ, ತನ್ನನ್ನು ತಾನೇ ಅರಿಯುವ ಹಾಗೂ ಸಕಲ ಸೃಷ್ಟಿಯನ್ನೂ ಪೋಷಿಸುವ ಆ ಉನ್ನತ ಶಕ್ತಿಯನ್ನು ತಿಳಿಯುವಲ್ಲಿ ನಿಜವಾದ ಆಸಕ್ತಿ ಬೆಳೆಯಬಹುದು. ಸರಿಯಾದ ಮಾರ್ಗದರ್ಶನದೊಂದಿಗೆ, ಅಂತಹ ಅನ್ವೇಷಕನು ಸರಿಯಾದ ವಿವೇಕವನ್ನು ಪಡೆದುಕೊಳ್ಳುತ್ತಾನೆ, ಇದು ಆಧ್ಯಾತ್ಮಿಕ ಹಾದಿಯಲ್ಲಿ ಅವರ ಪ್ರೇರಣೆಯನ್ನು ಮತ್ತಷ್ಟು ಪರಿಷ್ಕರಿಸುತ್ತದೆ. ಜೀವನದಲ್ಲಿ ನಮ್ಮ ಗುರಿಗಳು ಮತ್ತು ಆಸೆಗಳು, ನಮ್ಮ ತಿಳುವಳಿಕೆಯ ಮಟ್ಟವನ್ನು ಆಧರಿಸಿವೆ.
\end{inspiration}
\newpage

\slcol{\Index{ಬಹೂನಾಂ ಜನ್ಮನಾಮಂತೇ} ಜ್ಞಾನವಾನ್ಮಾಂ ಪ್ರಪದ್ಯತೇ ।\\
ವಾಸುದೇವಃ ಸರ್ವಮಿತಿ ಸ ಮಹಾತ್ಮಾ ಸುದುರ್ಲಭಃ ॥ ೧೯ ॥}
\cquote{ಅನೇಕ ಜನ್ಮಗಳ ಸಾಧನೆಯ ಕೊನೆಯಲ್ಲಿ ಜ್ಞಾನ ಪಡೆದು ಜ್ಞಾನಿಯು ಕೊನೆಗೆ ನನ್ನನ್ನು ಸೇರುವನು. ಎಲ್ಲಡೆಯೂ ಭಗವಂತನನ್ನೇ ಕಾಣುವ ಅಂತ ಮಹಾತ್ಮರು ತೀರ ವಿರಳ.}
\slcol{\Index{ಕಾಮೈಸ್ತೈಸ್ತೈರ್ಹೃತಜ್ಞಾನಾಃ} ಪ್ರಪದ್ಯಂತೇಽನ್ಯದೇವತಾಃ ।\\
ತಂ ತಂ ನಿಯಮಮಾಸ್ಥಾಯ ಪ್ರಕೃತ್ಯಾ ನಿಯತಾಃ ಸ್ವಯಾ ॥ ೨೦ ॥}
\cquote{ಬೇರೆ ಬೇರೆ ಬಯಕೆಗಳಿಗೆ ಬಲಿಯಾಗಿ ವಿವೇಕ ಕಳೆದುಕೊಂಡವರು, ತಮ್ಮ ತಮ್ಮ ಸ್ವಭಾವಕ್ಕನುಗುಣವಾಗಿ ಆಯಾ ಕಟ್ಟುಪಾಡುಗಳನ್ನು ಅವಲಂಬಿಸಿ ಬೇರೆ ಬೇರೆ ದೇವತೆಗಳನ್ನು ಭಜಿಸುತ್ತಾರೆ.}
\slcol{\Index{ಯೋ ಯೋ ಯಾಂ ಯಾಂ} ತನುಂ ಭಕ್ತಃ ಶ್ರದ್ಧಯಾರ್ಚಿತುಮಿಚ್ಛತಿ ।\\
ತಸ್ಯ ತಸ್ಯಾಚಲಾಂ ಶ್ರದ್ಧಾಂ ತಾಮೇವ ವಿದಧಾಮ್ಯಹಮ್ ॥ ೨೧ ॥}
\cquote{ಯಾವ ಭಕ್ತನು ಯಾವ ದೇವತೆಯನ್ನು ನಂಬಿ ಪೂಜಿಸುವನೋ ಅವನಿಗೆ ನಾನು ಆಯಾ ದೇವತೆಯ ಮೇಲಿನ ನಂಬಿಕೆಯನ್ನು ನೆಲೆಗೊಳಿಸುವೆನು.}
\slcol{\Index{ಸ ತಯಾ ಶ್ರದ್ಧಯಾ} ಯುಕ್ತಸ್ತಸ್ಯಾರಾಧನಮೀಹತೇ ।\\
ಲಭತೇ ಚ ತತಃ ಕಾಮಾನ್ಮಯೈವ ವಿಹಿತಾನ್ಹಿ ತಾನ್ ॥ ೨೨ ॥}
\cquote{ಅವನು ಆ ನಂಬಿಕೆಯಿಂದ ಕೂಡಿದ ಆ ದೇವತೆಯ ಪೂಜೆಯನ್ನು ಮಾಡುತ್ತಾನೆ, ನಾನು ಗೊತ್ತು ಮಾಡಿರುವಂತೆ ಆ ದೇವತೆಯಿಂದ ಫಲಗಳನ್ನೂ ಪಡೆಯುತ್ತಾನೆ.}
\slcol{\Index{ಅಂತವತ್ತು ಫಲಂ} ತೇಷಾಂ ತದ್ಭವತ್ಯಲ್ಪಮೇಧಸಾಮ್ ।\\
ದೇವಾಂದೇವಯಜೋ ಯಾಂತಿ ಮದ್ಭಕ್ತಾ ಯಾಂತಿ ಮಾಮಪಿ ॥ ೨೩ ॥}
\cquote{ದಡ್ಡರಾದ  ಅವರಿಗೆ ದೊರೆಯುವ ಫಲ ಮಾತ್ರ ಕ್ಷಣಿಕವಾದದ್ದು. ದೇವತೆಗಳನ್ನು ಪೂಜಿಸುವವರು ದೇವತೆಗಳನ್ನು ಪಡೆಯುತ್ತಾರೆ. ನನ್ನ ಭಕ್ತರು ನನ್ನನ್ನು ಪಡೆಯುತ್ತಾರೆ.}

\newpage
\begin{mananam}{\mananamTitle{ಮನನ ಶ್ಲೋಕ - ೧೮, ೧೯}}
\nfntmtxt ನನ್ನ ವಿವೇಕವನ್ನು  ಹೆಚ್ಚಿಸಿ, ಆಧ್ಯಾತ್ಮಿಕ ಹಾದಿಯಲ್ಲಿ ಸಾಕ್ಷಾತ್ಕಾರವನ್ನು ಪಡೆಯುವುದು ಹೇಗೆ? ಜ್ಞಾನಮಾರ್ಗದಲ್ಲಿ ಮುನ್ನಡೆಯುವಾಗ ಭಕ್ತಿ ಮತ್ತು ಏಕಾಗ್ರತೆಯ ನಡುವೆ ಸಮತೋಲನವನ್ನು ನಿರ್ವಹಿಸುವುದು ಏಕೆ ಮುಖ್ಯ? ‘ಜ್ಞಾನಮಾರ್ಗ’ವನ್ನು (ಜ್ಞಾನ ಯೋಗ) ಅದ್ವಿತೀಯಗೊಳಿಸುವುದು ಯಾವುದು?  ದೇವರು ಈ ಮಾರ್ಗದಲ್ಲಿ ಇರುವವರನ್ನು  ತನಗೆ ಅತೀ ಪ್ರೀತಿ ಪಾತ್ರರು ಮತ್ತು ಆತ್ಮೀಯರು ಎಂದು ಏಕೆ ಕೊಂಡಾಡುತ್ತಾನೆ? 
ಆಧ್ಯಾತ್ಮಿಕವಾಗಿ ಬುದ್ಧಿವಂತನಾದ ವ್ಯಕ್ತಿಯ ಯಾವ ದೂರದೃಷ್ಟಿ ಮತ್ತು ದೃಷ್ಟಿಕೋನವನ್ನು ಹೊಂದಲು ನಾನು ಆಶಿಸಬಹುದು? 
\end{mananam}
\sReflect
\begin{inspiration}{\mananamTitle{ಸ್ಪೂರ್ತಿ}}
\nfntmtxt ವಿವೇಕದ ಪರಾಕಾಷ್ಠೆಯು, ಆ ದೈವದೊಂದಿಗಿನ ಏಕತೆಯನ್ನು ಅರಿತುಕೊಳ್ಳುವುದು ಮತ್ತು ಎಲ್ಲದರಲ್ಲೂ ಆ ದೈವತ್ವವನ್ನು ಗ್ರಹಿಸುವುದೇ ಆಗಿದೆ. ಆಗ ಮಾತ್ರ  ಹುಡುಕಾಟಗಳೆಲ್ಲವೂ ಕೊನೆಗೊಳ್ಳುತ್ತವೆ.
ಮತ್ತೆಲ್ಲಾ ರೀತಿಯ ಪ್ರಾಪಂಚಿಕ ಜ್ಞಾನ ಮತ್ತು ಅವನ್ನು ಪೋಷಿಸುವ ಶಕ್ತಿಗಳು, ಇನ್ನೂ ದ್ವಂದ್ವದೃಷ್ಟಿಯನ್ನು ಸಮರ್ಥಿಸುತ್ತವೆ. ದ್ವಂದ್ವತೆ ಇನ್ನೂ ಇರುವಲ್ಲಿ, ಸಂಸಾರದ ದುಃಖವು ಮುಂದುವರಿಯುತ್ತದೆ. ಕೆಳಸ್ತರದ ಸತ್ಯಗಳು ಮತ್ತು ಅದರ ಪ್ರತಿಫಲಗಳೊಂದಿಗೆ ರಾಜಿ ಮಾಡಿಕೊಳ್ಳುವುದನ್ನು ನಿರಾಕರಿಸಿದಾಗ ಮಾತ್ರ, ಪರಮ ಜ್ಞಾನವನ್ನು ಸಾಧಿಸಲು ಸಾಧ್ಯ.
\end{inspiration}
\newpage

\begin{mananam}{\mananamTitle{ಮನನ ಶ್ಲೋಕ - ೨೧}}
\nfntmtxt ನನ್ನ ನಂಬಿಕೆಯ ಗುಣ-ಲಕ್ಷಣಗಳೇನು? ಇದು ಕೇವಲ ಭಾವನೆ, ನಿಷ್ಠೆಯ ಪ್ರಜ್ಞೆ ಅಥವಾ ಕುರುಡು ನಂಬಿಕೆಯೇ?  ನಾನು ಈ ಮಾರ್ಗವನ್ನು  ಅನುಸರಿಸಲು  ಹೇಗೆ ನಿರ್ಧರಿಸಿದೆ? ನಾನು ಅದರ  ಬಗ್ಗೆ ಸಾಕಷ್ಟು ಆಳವಾಗಿ ಚಿಂತನೆ ಮಾಡಿದ್ದೇನೆಯೇ?  ನಾನು, ನನ್ನಲ್ಲಿ ಮತ್ತು ನನ್ನ ದೈನಂದಿನ ಜೀವನದಲ್ಲಿ ಯಾವುದೇ ಪ್ರಗತಿ ಅಥವಾ, ಸಕಾರಾತ್ಮಕ ಬದಲಾವಣೆಗಳನ್ನು ಗಮನಿಸಿದ್ದೇನೆಯೇ? ನನ್ನ ಬುದ್ಧಿಶಕ್ತಿಯ ತಿಳುವಳಿಕೆ, ನನ್ನ ಅಂತಃಪ್ರಜ್ಞೆಯೊಂದಿಗೆ ಹೊಂದಿಕೆಯಾಗುತ್ತದೆಯೇ? ಈ ಮಾರ್ಗವು, ಸರ್ವೋಚ್ಚ ಸತ್ಯವನ್ನು ಮತ್ತು ಅದನ್ನು ಸಾಧಿಸುವ ವಿಧಾನಗಳ ದೃಷ್ಟಿಯನ್ನು ನೀಡುತ್ತದೆಯೇ? ಅಥವಾ, ಈ ಮಾರ್ಗವು ನನ್ನನ್ನು ಏಕಪಕ್ಷೀಯವಾಗಿ ಕಟ್ಟಿಹಾಕಿ ಇತರ ಮಾರ್ಗಗಳು ಮತ್ತು ವಿಧಾನಗಳನ್ನು ಕೀಳಾಗಿ ಕಾಣುವಂತೆ ಮಾಡುತ್ತಿದೆಯೇ?
\end{mananam}
\sReflect
\begin{inspiration}{\mananamTitle{ಸ್ಪೂರ್ತಿ}}
\nfntmtxt ಒಬ್ಬರ ತಿಳುವಳಿಕೆ ಗಾಢವಾಗುತ್ತಿದ್ದಂತೆ, ಅವರ ಲೌಕಿಕ ಮತ್ತು ಆಧ್ಯಾತ್ಮಿಕ ಜೀವನದಲ್ಲಿಯ ಆಕಾಂಕ್ಷೆಗಳೂ ಗಾಢವಾಗುತ್ತವೆ, ಯಾವ ದಾರಿಯಲ್ಲೇ  ಆಗಲಿ ನಮ್ಮ ಬದ್ಧತೆ ಆಳವಾದಂತೆಲ್ಲ ನಮ್ಮ ನಂಬಿಕೆಯು ತೀವ್ರವಾಗುತ್ತದೆ; ಇದರ ಪರಿಣಾಮವಾಗಿ, ಒಂದೇ ರೀತಿಯ ಸನಾತನ ಗ್ರಂಥಗಳನ್ನು ಅನುಸರಿಸುವವರಲ್ಲಿ ಕೂಡ  ವಿವಿಧ ‘ದರ್ಶನಗಳು’ (ದೃಷ್ಟಿಕೋನಗಳು ಅಥವಾ ತತ್ತ್ವಚಿಂತನೆಗಳು) ಇವೆ, ಪ್ರತಿಯೊಬ್ಬರೂ ತಾವೇ ಅಂತಿಮ ಸತ್ಯವನ್ನು ಹೊಂದಿದ್ದೇವೆಂದು  ಹೇಳಿಕೊಳ್ಳುತ್ತಾರೆ. ಈ ಶ್ಲೋಕವು, ಸರಿಯಾದ ಅಧ್ಯಯನ ಮತ್ತು ವಿವೇಚನೆಯಿಲ್ಲದೆ ನಮ್ಮ ನಂಬಿಕೆಯನ್ನು ತಪ್ಪು ವ್ಯಕ್ತಿಗಳ ಮತ್ತು ವಿಧಾನಗಳ ಮೇಲೆ ಇಡುವುದರ ವಿರುದ್ಧ ಎಚ್ಚರಿಕೆ ನೀಡುತ್ತದೆ. ಒಬ್ಬರ ಪಥದಲ್ಲಿ ನಿಷ್ಠರಾಗಿರುವುದು ಪ್ರಶಂಸಾರ್ಹವಾಗಿದ್ದರೂ, ಆ ಮಾರ್ಗವು ಸ್ವಮತಾಭಿಮಾನದಿಂದ ಇತರರ ಸಿದ್ಧಾಂತವನ್ನು ತಳ್ಳಿಹಾಕುವುದಾಗಿದ್ದರೆ ಅದು ಸಮಸ್ಯಾತ್ಮಕವಾಗುತ್ತದೆ; ಅಂತಹ ಸಂದರ್ಭಗಳಲ್ಲಿ, ನಾವು ನಮ್ಮ ದೃಷ್ಟಿಕೋಣವನ್ನು ಮರುಪರಿಶೀಲಿಸಬೇಕು.
\end{inspiration}
\newpage

\begin{mananam}{\mananamTitle{ಮನನ ಶ್ಲೋಕ - ೨೩}}
\nfntmtxt ಯಾವ ರೀತಿಯ ದೇವರು ಅಥವಾ ದೇವತೆಗಳನ್ನು ನಾನು ಪೂಜಿಸುತ್ತೇನೆ ಅಥವಾ ಪ್ರಾರ್ಥಿಸುತ್ತೇನೆ? ಈ ಅಭ್ಯಾಸಗಳಿಂದ ನನಗೆ ದೊರಕುವ ಫಲಿತಾಂಶವೇನು? ನನ್ನ ಅಭ್ಯಾಸಗಳು ಕೊಡುವ ಫಲಿತಾಂಶಗಳ ಭರವಸೆಯು, ಸೀಮಿತ ಹಾಗೂ ಸಮಯಾಧೀನವೇ? ನನ್ನ ಭಕ್ತಿಯು ಕೆಳ ಸ್ತರದ ದೇವತೆಗಳ ಕಡೆಗಿದೆಯೇ ಅಥವಾ, ಆ ಸರ್ವೋಚ್ಛ ಪರಮಾತ್ಮನೆಡೆಗಿದೆಯೇ? ಪರಮಾತ್ಮ, ಆತ್ಮ ಮತ್ತು ಪರಮ ಸತ್ಯದ ಬಗ್ಗೆ ನನಗೆ ಸಾಕಷ್ಟು ಅರಿವಿದೆಯೇ?
\end{mananam}
\sReflect
\begin{inspiration}{\mananamTitle{ಸ್ಪೂರ್ತಿ}}
\nfntmtxt ನಮ್ಮ ಅಭ್ಯಾಸಗಳ ಸ್ವರೂಪವನ್ನು ಲೆಕ್ಕಿಸದೇ, ನಮ್ಮ ಅಸ್ತಿತ್ವದ ಪರಮ ಸತ್ಯವಾದ ಆ ಪರಮಾತ್ಮನನ್ನು ನಾವು  ಅರಿತುಕೊಳ್ಳಲು ಆಶಿಸಬೇಕು. ಇತರ ಎಲ್ಲಾ ವಿಧದ ಪೂಜೆಗಳು ಮತ್ತು ಅಭ್ಯಾಸಗಳು ಅಲ್ಪಕಾಲದ ಪ್ರತಿಫಲವನ್ನು ಮಾತ್ರ ನೀಡುತ್ತವೆ.
ಪರಂಪರಾಯುಕ್ತ ಗುರುಗಳಿಂದ ಮತ್ತು ಧರ್ಮ ಗ್ರಂಥಗಳಿಂದ ಪಡೆದ ನಿಜವಾದ ಜ್ಞಾನವು ನಮ್ಮನ್ನು ಸರಿಯಾದ ದಾರಿಯಲ್ಲಿ ನಿಲ್ಲುವಂತೆ ಮತ್ತು ನಮ್ಮ ಪ್ರಯತ್ನಗಳು ಜೀವನದ ಸರ್ವೋಚ್ಚ ಗುರಿಯಿಂದ ದೂರವಾಗದಂತೆ ನೋಡಿಕೊಳ್ಳುತ್ತದೆ. 
\end{inspiration}
\newpage

\slcol{\Index{ಅವ್ಯಕ್ತಂ ವ್ಯಕ್ತಿಮಾಪನ್ನಂ} ಮನ್ಯಂತೇ ಮಾಮಬುದ್ಧಯಃ ।\\
ಪರಂ ಭಾವಮಜಾನಂತೋ ಮಮಾವ್ಯಯಮನುತ್ತಮಮ್ ॥ ೨೪ ॥}
\cquote{ನಾಶವಿರದ, ಹಿರಿದಾದ ನನ್ನ ನಿಜತತ್ವವನ್ನು ತಿಳಿಯದ ಅವಿವೇಕಿಗಳು ಅಪ್ರಾಕೃತಸ್ವರೂಪನಾದ ನನ್ನನ್ನು ಹುಟ್ಟು ಸಾವುಗಳಿಗೆ ಪಕ್ಕಾದ ಜಡದೇಹಿ ಎಂದು ತಿಳಿಯುತ್ತಾರೆ.}
\slcol{\Index{ನಾಹಂ ಪ್ರಕಾಶಃ ಸರ್ವಸ್ಯ} ಯೋಗಮಾಯಾಸಮಾವೃತಃ ।\\
ಮೂಢೋಽಯಂ ನಾಭಿಜಾನಾತಿ ಲೋಕೋ ಮಾಮಜಮವ್ಯಯಮ್ ॥ ೨೫ ॥}
\cquote{ನಿಜ ಸಾಮರ್ಥ್ಯದಿಂದ ಮತ್ತು ಮಾಯೆಯಿಂದ ಮರೆಯಾಗಿ ನಿಂತಿರುವ ನಾನು ಎಲ್ಲರಿಗೂ ಕಾಣಬರುವವನಲ್ಲ. ದಡ್ಡರಾದ ಈ ಜಗತ್ತಿನ ಜನರು ಹುಟ್ಟು ಸಾವುಗಳಿಲ್ಲದ ನನ್ನನ್ನು ತಿಳಿಯಲಾರರು.}
\slcol{\Index{ವೇದಾಹಂ ಸಮತೀತಾನಿ} ವರ್ತಮಾನಾನಿ ಚಾರ್ಜುನ ।\\
ಭವಿಷ್ಯಾಣಿ ಚ ಭೂತಾನಿ ಮಾಂ ತು ವೇದ ನ ಕಶ್ಚನ ॥ ೨೬ ॥}
\cquote{ಅರ್ಜುನ, ಹಿಂದಿನ, ಈಗಿನ, ಮುಂದಿನ ಎಲ್ಲವನ್ನೂ   ನಾನು ತಿಳಿದಿದ್ದೇನೆ. ನನ್ನನ್ನಾದರೋ ಯಾರೂ ಅರಿಯರು.}
\slcol{\Index{ಇಚ್ಛಾದ್ವೇಷಸಮುತ್ಥೇನ} ದ್ವಂದ್ವಮೋಹೇನ ಭಾರತ ।\\
ಸರ್ವಭೂತಾನಿ ಸಂಮೋಹಂ ಸರ್ಗೇ ಯಾಂತಿ ಪರಂತಪ ॥ ೨೭ ॥}
\cquote{ಹೇ ಪರಂತಪ, ಬೇಕು ಬೇಡ ಎಂಬ ಭಾವನೆಗಳಿಂದ ಹುಟ್ಟುವಂಥ ಸುಖ ದುಃಖ ಮೊದಲಾದ ದ್ವಂದ್ವದ ಆವೇಶಕ್ಕೂಳಗಾಗಿ ಎಲ್ಲಾ ಪ್ರಾಣಿಗಳೂ ಸೃಷ್ಟಿಯ ಆರಂಭದಲ್ಲಿಯೇ ಅಜ್ಞಾನದ ಸುಳಿಯಲ್ಲಿ ಬೀಳುತ್ತವೆ.}
\slcol{\Index{ಯೇಷಾಂ ತ್ವಂತಗತಂ} ಪಾಪಂ ಜನಾನಾಂ ಪುಣ್ಯಕರ್ಮಣಾಮ್ ।\\
ತೇ ದ್ವಂದ್ವಮೋಹನಿರ್ಮುಕ್ತಾ ಭಜಂತೇ ಮಾಂ ದೃಢವ್ರತಾಃ ॥ ೨೮ ॥}
\cquote{ಒಳ್ಳೆಯ ಕರ್ಮಗಳನ್ನು ಮಾಡುತ್ತಿರುವ ಯಾವ ಜನರ ಪಾಪಗಳೆಲ್ಲ ಕೊನೆಗೊಂಡಿರುವವೋ ಅವರು, ಆ ದ್ವಂದ್ವದ ಆವೇಶದಿಂದ ಬಿಡುಗಡೆ ಹೊಂದಿ ಕಟ್ಟುನಿಟ್ಟಾದ ವ್ರತವನ್ನು ಹಿಡಿದು ನನ್ನನ್ನು ಭಜಿಸುತ್ತಾರೆ.}
\newpage

\slcol{\Index{ಜರಾಮರಣಮೋಕ್ಷಾಯ} ಮಾಮಾಶ್ರಿತ್ಯ ಯತಂತಿ ಯೇ ।\\
ತೇ ಬ್ರಹ್ಮ ತದ್ವಿದುಃ ಕೃತ್ಸ್ನಮಧ್ಯಾತ್ಮಂ ಕರ್ಮ ಚಾಖಿಲಮ್ ॥ ೨೯ ॥}
\cquote{ಯಾರು ನನಗೆ ಶರಣು ಬಂದು ಮುಪ್ಪು ಸಾವುಗಳಿಂದ ಪಾರಾಗಬೇಕೆಂದು ಯತ್ನಿಸುವರೋ ಅವರು, ಆ ಪರಬ್ರಹ್ಮ ಸ್ವರೂಪನನ್ನೂ, ಆಧ್ಯಾತ್ಮವನ್ನೂ, ಎಲ್ಲ ಕರ್ಮಗಳನ್ನೂ ಸಂಪೂರ್ಣವಾಗಿ ತಿಳಿದುಕೊಳ್ಳುವರು.}
\slcol{\Index{ಸಾಧಿಭೂತಾಧಿದೈವಂ ಮಾಂ} ಸಾಧಿಯಙ್ಞಂ ಚ ಯೇ ವಿದುಃ ।\\
ಪ್ರಯಾಣಕಾಲೇಽಪಿ ಚ ಮಾಂ ತೇ ವಿದುರ್ಯುಕ್ತಚೇತಸಃ ॥ ೩೦ ॥}
\cquote{ಅಧಿಭೂತ, ಅಧಿದೈವ, ಅಧಿಯಜ್ಞಗಳೊಡಗೂಡಿದ ನನ್ನನ್ನು ಯಾರು ತಿಳಿಯುತ್ತಾರೋ ನನ್ನಲ್ಲಿಯೇ ಮನಸ್ಸಿಟ್ಟಿರುವ ಅವರು, ಸಾವಿನ ಸಮಯದಲ್ಲಿ ಕೂಡ ನನ್ನನ್ನು ನೆನಪಿಡುವರು.}

\newpage
\begin{mananam}{\mananamTitle{ಮನನ ಶ್ಲೋಕ - ೨೯}}
\nfntmtxt ಈ ಪ್ರಪಂಚದಲ್ಲಿ ನನ್ನ ಬೆಂಬಲದ ಮೂಲಗಳು ಯಾರು ಮತ್ತು ಏನು? ನನ್ನ ಬೆಂಬಲಕ್ಕಾಗಿ ಯಾವ ಕ್ಷೇತ್ರಗಳಲ್ಲಿ ನಾನು ನನ್ನ ಕುಟುಂಬ, ಸ್ನೇಹಿತರು ಮತ್ತು ಇತರರನ್ನು ಅವಲಂಬಿಸಬಹುದು? ಯಾವ ಕ್ಷೇತ್ರಗಳಲ್ಲಿ  ಅವರನ್ನು ಅವಲಂಬಿಸಲು ನನಗೆ ಸಾಧ್ಯವಾಗಲಾರದು?  ಯಾವ ಪರಿಸ್ಥಿತಿಗಳಲ್ಲಿ ಲೌಕಿಕ ಆಶ್ರಯದ ಮೂಲಗಳು ಅಸಮರ್ಪಕ ಅಥವಾ ಸೀಮಿತವಾಗಿರುತ್ತವೆ?; ಈ ಮಿತಿಗಳನ್ನು ಮೀರಿದ ಆಶ್ರಯಕ್ಕಾಗಿ ನಾನು ಯಾವ ಪರಿಕಲ್ಪನೆಯನ್ನು ರೂಪಿಸಿಕೊಳ್ಳಬಹುದು? ಹೇಗೆ ಅಂತಹ ‘ಅನಂತ’ವು, ನನ್ನ ಆತ್ಮಜ್ಞಾನವನ್ನು ಅಭಿವೃದ್ಧಿಗೊಳಿಸುತ್ತಿರುವ ಅಂತಿಮ ಭರವಸೆಯಾಗಿರುವುದು? 
\end{mananam}
\sReflect
\begin{inspiration}{\mananamTitle{ಸ್ಪೂರ್ತಿ}}
\nfntmtxt ಪ್ರಾಪಂಚಿಕ ಜೀವನದ ಹಲವು ಹಂತಗಳಲ್ಲಿ ನಾವು, ಜನ, ಸಮಾಜ ಮತ್ತು ಸರ್ಕಾರದ ಬೆಂಬಲವನ್ನು ಅವಲಂಬಿಸಬಹುದು. ಆದಾಗ್ಯೂ, ವೃದ್ಧಾಪ್ಯ ಮತ್ತು ಸಾವಿನ ಅನಿವಾರ್ಯತೆಯನ್ನು ಎದುರಿಸುವಾಗ, ನಮ್ಮ ಅತ್ಯುನ್ನತ ಆಶ್ರಯವು ದೇವರು ಮಾತ್ರ. ನಾವು ಎಷ್ಟೇ ತರ್ಕಬದ್ಧವಾಗಿ ವಿಚಾರ ಮಾಡಿದಾಗ್ಯೂ, ಯಾವುದಾದರೂ ವಿಷಯಗಳನ್ನು ಹಿಡಿದಿಟ್ಟುಕೊಳ್ಳುಲೇಬೇಕೆಂದು ನಮ್ಮ ಮನಸ್ಸಿಗೆ ಎನಿಸುತ್ತದೆ.  ಒಬ್ಬ ಸಾಧಕನಿಗೆ, ದೇವರ ಸ್ವರೂಪ ಅಥವಾ ದೈವತ್ವದ ಇರುವಿಕೆಯ ಭಾವನೆಯು ಆ ಪ್ರೋತ್ಸಾಹವನ್ನು  ನೀಡುತ್ತದೆ. ವಿವೇಕದಲ್ಲಿ ದೃಢವಾಗಿರುವವನಿಗೆ, ಅವನ ಆತ್ಮಜ್ಞಾನ ಮತ್ತು ಅದರ ಅನಂತನೊಂದಿಗಿನ ಭಗ್ನವಾಗದ ಸಂಪರ್ಕವು, ಒಂದು ಭದ್ರವಾದ ಮೂಲದಂತೆ ಕಾರ್ಯನಿರ್ವಹಿಸುತ್ತದೆ.
\end{inspiration}
\newpage

\chapEndSloka{ಜ್ಞಾನ ವಿಜ್ಞಾನಯೋಗ}

